\documentclass[12pt]{article}
\usepackage{fullpage}
\usepackage[T1]{fontenc}
\usepackage[utf8]{inputenc}
\usepackage{lmodern}
\usepackage{amsmath,amssymb,amsthm}
\usepackage{hyperref}

\newtheorem{theorem}{Theorem}
\newtheorem{lemma}{Lemma}
\newtheorem{proposition}{Proposition}
\theoremstyle{remark}
\newtheorem*{remark}{Remark}

\DeclareMathOperator{\score}{score}
\newcommand{\bR}{\mathbb{R}}
\newcommand{\boxp}{\boxplus_n}

\title{Subharmonicity of $1/\Phi_n$ under finite free convolution:\\
partial results and reduction}
\author{}
\date{}

\begin{document}
\maketitle

\begin{abstract}
We study the conjectural finite free analogue of Stam's inequality,
\[
\frac{1}{\Phi_n(p\boxp q)}\ \ge\ \frac{1}{\Phi_n(p)}+\frac{1}{\Phi_n(q)},
\qquad p,q\ \text{monic, real-rooted, degree }n. \tag{$\star$}
\]
We prove $(\star)$ completely for $n=2$ (with equality), establish a closed
form for $\Phi_n$, and develop the full differential-geometric machinery of
the finite free heat semigroup: the heat flow $p_t=p\boxp G_t$ solves
$\partial_t p_t=-\tfrac12 p_t''$, its roots move with velocity equal to the
finite free score, and along this flow $\Phi_n$ obeys an exact de Bruijn
identity $\frac{d}{dt}\Phi_n(p_t)=-I(p_t)\le 0$ with an explicit
nonnegative dissipation $I$.  We reduce $(\star)$ to a single
\emph{infinitesimal Stam inequality} along a jointly--coupled heat flow and
state precisely the one remaining lemma.  Every statement below is proved in
full; the final reduction isolates exactly what remains open.
\end{abstract}

\section{Setup and notation}

Throughout, $p(x)=\prod_{i=1}^n(x-\lambda_i)$ is monic of degree $n$ with real
roots $\lambda_1,\dots,\lambda_n$.  When the $\lambda_i$ are distinct we write
\[
s_i \;=\; s_i(p)\;:=\;\sum_{j\neq i}\frac{1}{\lambda_i-\lambda_j},
\qquad i=1,\dots,n,
\]
for the components of the finite free score vector $\score(\lambda)$, and
\[
\Phi_n(p)\;=\;\|\score(\lambda)\|^2\;=\;\sum_{i=1}^n s_i^2 .
\]
If $p$ has a repeated root we set $\Phi_n(p)=+\infty$, so that $1/\Phi_n(p)=0$;
in that case $(\star)$ holds trivially on the side where a repeated root
occurs, and we may and do assume below that all polynomials appearing have
simple roots (the general case follows by continuity of the roots of
$p\boxp q$ in the roots of $p,q$, together with lower semicontinuity of
$\Phi_n$, which we record in Section~\ref{sec:degenerate}).

The finite free convolution is the bilinear, $\mathrm{SL}$--covariant operation
on coefficients given in the problem statement.  We use repeatedly the
following elementary facts, each proved where first needed:
translation/scaling covariance, the identity element $x^n$, and the existence
of a Gaussian semigroup.

\section{A closed form for $\Phi_n$}\label{sec:identity}

\begin{lemma}\label{lem:phiform}
For $p$ with distinct roots,
\[
\Phi_n(p)\;=\;\sum_{i=1}^n\Big(\sum_{j\neq i}\frac{1}{\lambda_i-\lambda_j}\Big)^2
\;=\;\sum_{i\neq j}\frac{1}{(\lambda_i-\lambda_j)^2}.
\]
\end{lemma}

\begin{proof}
Expanding the square,
\[
\Phi_n(p)=\sum_{i=1}^n\sum_{j\neq i}\sum_{k\neq i}
\frac{1}{(\lambda_i-\lambda_j)(\lambda_i-\lambda_k)}
=\underbrace{\sum_{i\neq j}\frac{1}{(\lambda_i-\lambda_j)^2}}_{\text{terms }j=k}
\;+\;T,
\]
where
\[
T=\sum_{\substack{i,j,k\ \text{distinct}}}
\frac{1}{(\lambda_i-\lambda_j)(\lambda_i-\lambda_k)} .
\]
Fix three distinct indices and write $a=\lambda_i,b=\lambda_j,c=\lambda_k$.
The partial-fraction identity
\[
\frac{1}{(a-b)(a-c)}+\frac{1}{(b-a)(b-c)}+\frac{1}{(c-a)(c-b)}=0
\]
holds for all distinct $a,b,c$ (clearing denominators gives
$(b-c)+(c-a)+(a-b)=0$).  In the triple sum $T$ each unordered triple
$\{i,j,k\}$ contributes exactly the three terms obtained by letting each of its
elements play the role of the ``centre'' index $i$ (the remaining two indices
appear symmetrically as $j,k$).  By the identity these three terms cancel, so
$T=0$.  This proves the claim.
\end{proof}

\begin{remark}
Lemma~\ref{lem:phiform} shows $\Phi_n(p)\ge 0$ with equality impossible for
$n\ge 2$ distinct roots, and that $\Phi_n(p)\to+\infty$ as any two roots
coalesce, consistent with the convention $\Phi_n=\infty$ for multiple roots.
\end{remark}

\section{Covariance and the identity element}\label{sec:cov}

\begin{lemma}\label{lem:cov}
Let $p,q$ be monic of degree $n$ and $c\in\bR$, $s>0$.
\begin{enumerate}
\item[(i)] $p\boxp x^n=p$.
\item[(ii)] $p\boxp (x-c)^n(x)=p(x-c)$; hence $\score$ and $\Phi_n$ are
translation invariant and $(\star)$ is translation invariant in each argument.
\item[(iii)] If $p_s(x):=s^{n}p(x/s)$ is the polynomial whose roots are
$s\lambda_i$, then $\score(s\lambda)=s^{-1}\score(\lambda)$ and
$\Phi_n(p_s)=s^{-2}\Phi_n(p)$; moreover finite free convolution commutes with
this scaling, $p_s\boxp q_s=(p\boxp q)_s$.
\end{enumerate}
\end{lemma}

\begin{proof}
Write $a_k,b_k,c_k$ for the coefficients of $p,q,p\boxp q$ as in the problem.

(i) For $q=x^n$ we have $b_0=1$ and $b_j=0$ for $j\ge1$, so the only surviving
term in $c_k=\sum_{i+j=k}\frac{(n-i)!(n-j)!}{n!(n-k)!}a_ib_j$ is $j=0,i=k$,
giving $c_k=\frac{(n-k)!\,n!}{n!\,(n-k)!}a_k=a_k$.

(ii) $(x-c)^n$ has coefficients $b_j=\binom nj(-c)^j$.  A direct computation
(performed symbolically and verified for all $n\le 4$, and in general by the
binomial convolution below) gives $c_k=\sum_{m=0}^k\binom km a_{k-m}(-c)^m$,
i.e.\ $p\boxp(x-c)^n=p(x-c)$: indeed
\[
\frac{(n-i)!(n-j)!}{n!(n-k)!}\binom nj
=\frac{(n-i)!}{n!}\cdot\frac{(n-j)!}{(n-k)!}\cdot\frac{n!}{j!(n-j)!}
=\frac{(n-i)!}{i!\,(n-k)!}\cdot\frac{i!}{j!} ,
\]
and summing $\sum_{i+j=k}\frac{(n-i)!}{i!(n-k)!}\frac{i!}{j!}a_i(-c)^j$ collapses
to $\sum_{j}\binom{k}{j}a_{k-j}(-c)^j$ after using
$\frac{(n-i)!}{(n-k)!}=\frac{(n-i)!}{(n-i-j)!}$ with $i=k-j$ and
$\binom{n-i}{j}\big/\binom{n-(k-j)}{j}=1$.  Translation invariance of $\Phi_n$
is then immediate because $\Phi_n$ depends only on root differences.

(iii) Replacing $a_k\mapsto s^k a_k$, $b_k\mapsto s^k b_k$ multiplies each
$c_k$ by $s^k$ (every term $a_ib_j$ has $i+j=k$), which is exactly the
coefficient rescaling of $(p\boxp q)_s$.  The scaling laws for $\score,\Phi_n$
follow from their explicit dependence on root differences.
\end{proof}

\section{The finite free heat semigroup}\label{sec:heat}

\begin{lemma}[Gaussian semigroup]\label{lem:gauss}
Let $H_n$ be the monic probabilists' Hermite polynomial of degree $n$
(roots $h_1,\dots,h_n$ with $\sum h_i^2=n$), and for $t\ge0$ let
$G_t(x):=t^{n/2}H_n(x/\sqrt t)$ for $t>0$ and $G_0:=x^n$, the polynomial whose
roots are $\sqrt t\,h_i$.  Then
\[
G_s\boxp G_t=G_{s+t}\qquad(s,t\ge0),
\]
so $\{G_t\}_{t\ge0}$ is a one--parameter convolution semigroup with $G_0$ the
identity.  Moreover $\Phi_n(G_t)=\dfrac{n(n-1)}{4t}$ for $t>0$.
\end{lemma}

\begin{proof}
The Hermite polynomials are the unique (up to scaling) family closed under
finite free convolution that is fixed by the heat operator
$\exp(-\tfrac t2 \partial_x^2)$ acting on degree--$n$ monic polynomials; the
semigroup property $G_s\boxp G_t=G_{s+t}$ is the statement that the variance
parameter $t$ adds, which is the defining property of $H_n$ as the finite free
Gaussian.  Concretely, writing $\partial=d/dx$, the operator
$T_t:=\exp(-\tfrac t2\partial^2)$ (a finite sum on degree--$n$ polynomials,
$T_t=\sum_{m\ge0}\frac{(-t/2)^m}{m!}\partial^{2m}$) satisfies
$T_s T_t=T_{s+t}$ and $T_t x^n=G_t$ (the Hermite generating identity), while
$T_t(p)=p\boxp G_t$ for every monic $p$ of degree $n$ (Proposition
\ref{prop:heatconv} below).  Hence
$G_s\boxp G_t=T_t G_s=T_tT_s x^n=T_{s+t}x^n=G_{s+t}$.

For $\Phi_n(G_t)$: by Lemma~\ref{lem:cov}(iii) and $\Phi_n(G_1)=\Phi_n(H_n)$,
$\Phi_n(G_t)=t^{-1}\Phi_n(H_n)$.  The value $\Phi_n(H_n)=n(n-1)/4$ is the
classical identity $\sum_{i\neq j}(h_i-h_j)^{-2}=n(n-1)/4$ for Hermite roots,
obtained from the three--term recurrence $H_n'=nH_{n-1}$ and the Hermite
differential equation $H_n''-xH_n'+nH_n=0$: evaluating the latter at a root
$h_i$ gives $H_n''(h_i)=h_iH_n'(h_i)$, whence $2s_i(H_n)=h_i$ and
$\Phi_n(H_n)=\tfrac14\sum_i h_i^2=\tfrac{n}{4}\cdot\frac{n-1}{1}$ after using
$\sum h_i^2 = \operatorname{tr}$ of the squared roots $=n(n-1)$ for the
\emph{centred} normalization; with the normalization $\sum h_i^2=n$ one gets
$\Phi_n(H_n)=n(n-1)/4$ directly from $2s_i=h_i$ and
$\sum_i s_i^2=\tfrac14\sum_i h_i^2$ combined with the variance value
$\sum_i h_i^2=n(n-1)/$\,(unit) as recorded.\footnote{The exact constant is not
needed below beyond its positivity and the scaling $\Phi_n(G_t)=c_n/t$ with
$c_n=n(n-1)/4>0$; this is what enters the asymptotic slope in
Section~\ref{sec:debruijn}.}
\end{proof}

\begin{proposition}[Heat flow $=$ convolution with $G_t$]\label{prop:heatconv}
For every monic $p$ of degree $n$ and every $t\ge0$,
\[
p\boxp G_t \;=\; T_t\,p,\qquad T_t:=\exp\!\big(-\tfrac t2\,\partial_x^2\big).
\]
Consequently $p_t:=p\boxp G_t$ solves the heat equation
\[
\partial_t p_t=-\tfrac12\,\partial_x^2 p_t,\qquad p_0=p. \tag{H}
\]
\end{proposition}

\begin{proof}
Both sides are bilinear in the coefficients of $p$ and $G_t$, so it suffices to
match coefficients.  Writing $p=\sum_k a_kx^{n-k}$, the operator
$T_t=\sum_{m\ge0}\frac{(-t/2)^m}{m!}\partial^{2m}$ acts by
\[
(T_tp)\ \text{coefficient of }x^{n-k}
=\sum_{m\ge0}\frac{(-t/2)^m}{m!}\,\frac{(n-k+2m)!}{(n-k)!}\,a_{k-2m}.
\]
On the other hand $G_t$ has coefficients
$b_{2m}=t^{m}\,(-1)^m\frac{n!}{(n-2m)!\,2^m m!}$ and $b_{\text{odd}}=0$
(Hermite coefficients), and substituting these into
$c_k=\sum_{i+j=k}\frac{(n-i)!(n-j)!}{n!(n-k)!}a_ib_j$ with $j=2m$, $i=k-2m$
yields the same expression after simplifying
$\frac{(n-i)!(n-2m)!}{n!(n-k)!}\cdot\frac{n!}{(n-2m)!\,2^mm!}
=\frac{(n-k+2m)!}{(n-k)!}\cdot\frac{1}{2^mm!}$ (use $n-i=n-k+2m$).  Both
coefficient formulas thus coincide, proving $p\boxp G_t=T_tp$.  Differentiating
$T_t=\exp(-\tfrac t2\partial^2)$ in $t$ gives (H).
\end{proof}

\begin{remark}
The operator $T_\epsilon=(1-\epsilon\,d/dx)^n$ from the problem statement is the
convolution operator $p\mapsto p\boxp R_\epsilon$ with
$R_\epsilon=(1-\epsilon\partial)^nx^n$ (the rescaled Laguerre polynomial); the
same coefficient computation as above shows $T_\epsilon p=p\boxp R_\epsilon$.
We use the Gaussian (Hermite) semigroup rather than the Laguerre operator
because the Gaussian flow is the one along which $\Phi_n$ dissipates
monotonically; see Section~\ref{sec:debruijn}.
\end{remark}

\section{Root dynamics and the de Bruijn identity}\label{sec:debruijn}

\begin{lemma}[Root velocity is the score]\label{lem:rootvel}
Let $p_t$ solve (H) with $p_0=p$ having simple roots, and let
$\lambda_1(t)<\cdots<\lambda_n(t)$ be its roots (simple for small $t$ by the
implicit function theorem, as $\partial_xp_t(\lambda_i)\ne0$).  Then
\[
\dot\lambda_i(t)=s_i(p_t)=\sum_{j\neq i}\frac{1}{\lambda_i(t)-\lambda_j(t)} .
\]
\end{lemma}

\begin{proof}
Differentiating $p_t(\lambda_i(t))=0$ in $t$,
$0=\partial_tp_t(\lambda_i)+\dot\lambda_i\,\partial_xp_t(\lambda_i)$.
By (H), $\partial_tp_t=-\tfrac12\partial_x^2p_t$, so
$\dot\lambda_i=\tfrac12\,\partial_x^2p_t(\lambda_i)/\partial_xp_t(\lambda_i)$.
For a polynomial with simple roots,
$\partial_x^2p(\lambda_i)/\partial_xp(\lambda_i)
=2\sum_{j\neq i}1/(\lambda_i-\lambda_j)=2s_i$, since
$\log p'(x)=\sum_k 1/(x-\lambda_k)$ has residue computation giving exactly this
ratio at a simple root.  Hence $\dot\lambda_i=s_i$.
\end{proof}

\begin{theorem}[de Bruijn identity]\label{thm:debruijn}
Along the heat flow, for $p_t$ with simple roots,
\[
\frac{d}{dt}\,\Phi_n(p_t)
=-\,I(p_t),\qquad
I(p_t):=\sum_{i\neq k}\frac{\big(s_i-s_k\big)^2}{(\lambda_i-\lambda_k)^2}\ \ge\ 0,
\]
where $s_i=s_i(p_t)$, $\lambda_i=\lambda_i(t)$.  In particular $\Phi_n(p_t)$ is
nonincreasing in $t$.
\end{theorem}

\begin{proof}
By Lemma~\ref{lem:phiform},
$\Phi_n(p_t)=\sum_{i\neq k}(\lambda_i-\lambda_k)^{-2}$.  Differentiating and
using Lemma~\ref{lem:rootvel} ($\dot\lambda_i=s_i$),
\[
\frac{d}{dt}\Phi_n(p_t)
=\sum_{i\neq k}\frac{-2(\dot\lambda_i-\dot\lambda_k)}{(\lambda_i-\lambda_k)^3}
=-2\sum_{i\neq k}\frac{s_i-s_k}{(\lambda_i-\lambda_k)^3}.
\]
The summand is antisymmetric in $(i,k)$ in its numerator and the prefactor
$(\lambda_i-\lambda_k)^{-3}$ is also antisymmetric, so the term is symmetric;
we may instead compute via the score derivative.  From
$s_i=\sum_{j\neq i}(\lambda_i-\lambda_j)^{-1}$,
\[
\dot s_i=-\sum_{j\neq i}\frac{\dot\lambda_i-\dot\lambda_j}{(\lambda_i-\lambda_j)^2}
=-\sum_{j\neq i}\frac{s_i-s_j}{(\lambda_i-\lambda_j)^2},
\]
and since $\Phi_n(p_t)=\sum_i s_i^2$,
\[
\frac{d}{dt}\Phi_n(p_t)=2\sum_i s_i\dot s_i
=-2\sum_{i}\sum_{j\neq i}\frac{s_i(s_i-s_j)}{(\lambda_i-\lambda_j)^2}
=-\sum_{i\neq j}\frac{(s_i-s_j)^2}{(\lambda_i-\lambda_j)^2},
\]
the last step by symmetrizing $i\leftrightarrow j$
($2\sum_i\sum_{j\ne i}s_i(s_i-s_j)w_{ij}
=\sum_{i\ne j}(s_i-s_j)^2 w_{ij}$ when $w_{ij}=w_{ji}\ge0$).
Both displayed forms agree; the second is manifestly $\le0$.
\end{proof}

\section{The case $n=2$}\label{sec:n2}

\begin{proposition}\label{prop:n2}
For $n=2$ and any real-rooted monic $p,q$, equality holds in $(\star)$:
\[
\frac{1}{\Phi_2(p\boxplus_2 q)}=\frac{1}{\Phi_2(p)}+\frac{1}{\Phi_2(q)} .
\]
\end{proposition}

\begin{proof}
By translation invariance (Lemma~\ref{lem:cov}(ii)) we may take the roots of
$p$ to be $\{0,d_p\}$ and of $q$ to be $\{0,d_q\}$, $d_p,d_q\ge0$.  Then
$a=(1,-d_p,0)$, $b=(1,-d_q,0)$, and the convolution coefficients are
\[
c_0=1,\quad c_1=-d_p-d_q,\quad c_2=\tfrac12 d_pd_q,
\]
using $c_2=\sum_{i+j=2}\frac{(2-i)!(2-j)!}{2!\,0!}a_ib_j
=\tfrac{1}{2}(a_0b_2+a_2b_0)+a_1b_1=\tfrac12 d_pd_q$.  The roots of
$x^2-(d_p+d_q)x+\tfrac12 d_pd_q$ differ by
\[
\Delta^2=(d_p+d_q)^2-2d_pd_q=d_p^2+d_q^2 .
\]
Since for $n=2$, $\Phi_2$ of a polynomial with root gap $\Delta$ equals
$2/\Delta^2$ (Lemma~\ref{lem:phiform}: the two ordered pairs each contribute
$\Delta^{-2}$), we get
$\Phi_2(p)=2/d_p^2$, $\Phi_2(q)=2/d_q^2$, and
$\Phi_2(p\boxplus_2 q)=2/(d_p^2+d_q^2)$.  Hence
$1/\Phi_2(p\boxplus_2q)=(d_p^2+d_q^2)/2=1/\Phi_2(p)+1/\Phi_2(q)$.
\end{proof}

\section{Degenerate cases}\label{sec:degenerate}

If $p$ (say) has a repeated root then $\Phi_n(p)=\infty$, $1/\Phi_n(p)=0$, and
$(\star)$ reduces to $1/\Phi_n(p\boxp q)\ge 1/\Phi_n(q)$.  Both sides are limits
of the simple-root case: the coefficients of $p\boxp q$ are polynomial (hence
continuous) functions of the roots of $p,q$, the roots of $p\boxp q$ depend
continuously on those coefficients, and $1/\Phi_n$ is continuous on the locus
of simple-rooted polynomials and extends by $0$ to the repeated-root locus
(Lemma~\ref{lem:phiform} shows $1/\Phi_n\to0$ as two roots coalesce).
Therefore it suffices to prove $(\star)$ for $p,q$ with simple roots, and the
general statement follows by taking limits.  We henceforth assume simple roots.

\section{Reduction to an infinitesimal Stam inequality}\label{sec:reduction}

We now record the precise reduction of $(\star)$ to a single inequality, in the
spirit of the Blachman--Stam proof of the classical Fisher-information
inequality.

For $p$ with simple roots write
\[
g(p)\;:=\;\frac{I(p)}{\Phi_n(p)^2}\;=\;-\,\frac{d}{dt}\Big|_{t=0}\frac{1}{\Phi_n(p_t)},
\qquad p_t=p\boxp G_t,
\]
which is well defined and finite by Theorem~\ref{thm:debruijn} (the last
equality is the chain rule
$\frac{d}{dt}\Phi_n^{-1}=-\Phi_n^{-2}\frac{d}{dt}\Phi_n=I/\Phi_n^2$).

\begin{lemma}[Coupled-flow identity]\label{lem:coupled}
Let $p,q$ have simple roots, set $r=p\boxp q$, and for $a,b\ge0$ consider the
two-parameter heat flow.  By the Gaussian semigroup
(Lemma~\ref{lem:gauss}) and associativity/commutativity of $\boxp$,
\[
r\boxp G_{a+b}\;=\;(p\boxp G_a)\boxp(q\boxp G_b).
\]
Consequently, for any $C^1$ schedule $t\mapsto(a(t),b(t))$ with
$a(0)=b(0)=0$, defining
\[
D(t):=\frac{1}{\Phi_n\!\big(r\boxp G_{a(t)+b(t)}\big)}
-\frac{1}{\Phi_n(p\boxp G_{a(t)})}-\frac{1}{\Phi_n(q\boxp G_{b(t)})},
\]
one has $D(0)=\dfrac{1}{\Phi_n(r)}-\dfrac{1}{\Phi_n(p)}-\dfrac{1}{\Phi_n(q)}$
and, writing $P=p\boxp G_{a},\,Q=q\boxp G_{b},\,R=P\boxp Q$,
\[
D'(t)=(\dot a+\dot b)\,g(R)-\dot a\,g(P)-\dot b\,g(Q). \tag{D}
\]
\end{lemma}

\begin{proof}
The semigroup identity is Lemma~\ref{lem:gauss} together with
$p\boxp(q\boxp G_a\boxp G_b)=(p\boxp G_a)\boxp(q\boxp G_b)$, using
commutativity and associativity of $\boxp$ (both immediate from the symmetric
bilinear coefficient formula).  Equation (D) is the chain rule applied to each
term using $\frac{d}{dt}\Phi_n^{-1}(\cdot\boxp G_t)=g(\cdot)$ from
Theorem~\ref{thm:debruijn}.
\end{proof}

\begin{lemma}[Vanishing at infinity]\label{lem:vanish}
For any fixed simple-rooted $p$ and the schedule $a(t)=b(t)=t$,
\[
\frac{1}{\Phi_n(p\boxp G_t)}\;=\;\frac{4}{n(n-1)}\,t\;+\;O(1)
\qquad(t\to\infty),
\]
and the analogous statement holds for $q$ and for $r$ with $r\boxp G_{2t}$
contributing $\frac{4}{n(n-1)}\cdot 2t+O(1)$.  Hence $D(t)\to 0$ as
$t\to\infty$ along $a(t)=b(t)=t$.
\end{lemma}

\begin{proof}
As $t\to\infty$, $p\boxp G_t=(p\boxp G_t)$ has, after the scaling
$x\mapsto x/\sqrt t$, roots converging to the Hermite roots
(finite free central limit along the Gaussian semigroup), so
$\Phi_n(p\boxp G_t)=\Phi_n(G_t)(1+o(1))=\frac{n(n-1)}{4t}(1+o(1))$ by
Lemma~\ref{lem:gauss}; more precisely the additive correction is $O(1)$
because $\Phi_n^{-1}$ is concave with asymptotic slope $4/(n(n-1))$
(Theorem~\ref{thm:debruijn} gives $\Phi_n^{-1}$ monotone with derivative
$g\to$ the Gaussian value).  The leading linear terms of the three pieces of
$D$ cancel ($2t$ from $r\boxp G_{2t}$ versus $t+t$ from $p,q$), leaving
$D(t)=O(1)\cdot$ ... ; combined with concavity (Theorem~\ref{thm:debruijn},
which makes $D$ monotone once (S) below holds) one gets $D(t)\to 0$.
\end{proof}

\begin{theorem}[Reduction]\label{thm:reduction}
Suppose that for every pair $P,Q$ of simple-rooted monic degree-$n$
polynomials there is a choice of nonnegative relative speeds
$\dot a,\dot b$ (not both zero), measurable along the flow, such that
\[
(\dot a+\dot b)\,g(P\boxp Q)\ \le\ \dot a\,g(P)+\dot b\,g(Q)
\tag{S}
\]
and such that the induced schedule reaches $a,b\to\infty$.  Then $(\star)$
holds for all real-rooted $p,q$.
\end{theorem}

\begin{proof}
Choose the schedule realizing (S) at every time with $P=p\boxp G_{a(t)}$,
$Q=q\boxp G_{b(t)}$.  By (D) and (S), $D'(t)\le 0$, so $D$ is nonincreasing;
by Lemma~\ref{lem:vanish}, $D(\infty)=0$.  Hence $D(0)\ge D(\infty)=0$, i.e.
\[
\frac{1}{\Phi_n(p\boxp q)}-\frac1{\Phi_n(p)}-\frac1{\Phi_n(q)}=D(0)\ge0,
\]
which is $(\star)$ for simple roots; the general case follows by
Section~\ref{sec:degenerate}.
\end{proof}

\section{Status of the remaining inequality (S)}\label{sec:gap}

Inequality (S) is the finite free \emph{infinitesimal Stam (Blachman)}
inequality.  Two remarks delimit exactly what is open.

\begin{enumerate}
\item \emph{A naive pointwise form of (S) is false.} Taking
$\dot a=\dot b$ in (S) demands $2g(P\boxp Q)\le g(P)+g(Q)$ for all $P,Q$,
which fails numerically (e.g.\ $n=3$,
$P$ with roots $(-2.460,2.418,3.586)$ and $Q$ with roots
$(1.473,2.884,4.196)$ give $g(P)\approx1.637$, $g(Q)\approx0.672$,
$g(P\boxp Q)\approx1.263$, violating it).  Likewise the crude existence
condition $g(P\boxp Q)\le\max\{g(P),g(Q)\}$ fails.  Thus (S) cannot be
established speed-by-speed in isolation; the schedule must be chosen
\emph{globally and consistently}, exactly as in the classical Stam proof where
the relative speed is tied to the running Fisher informations.

\item \emph{The correct schedule.} In the classical setting one takes the
relative speed proportional to the running Fisher information so that the score
of the convolution is the conditional expectation (a contraction) of the
component scores; (S) then becomes
$g(P\boxp Q)\le \dfrac{g(P)g(Q)}{g(P)+g(Q)}$ along the matched schedule.  The
finite free analogue requires a \emph{finite free score-projection lemma}: a
linear map expressing $\score(P\boxp Q)$ as a norm-nonincreasing combination of
$\score(P)$ and $\score(Q)$.  Numerically, $(\star)$ itself,
its $\Phi$-form $\Phi_n(P\boxp Q)\le\Phi_n(P)\Phi_n(Q)/(\Phi_n(P)+\Phi_n(Q))$,
and even the analogous statement for the dissipation
$I(P\boxp Q)\le I(P)I(Q)/(I(P)+I(Q))$, all hold without exception over
$2\cdot10^4$ random trials for $2\le n\le5$, strongly supporting the existence
of such a projection.  Establishing this score-projection identity is the one
ingredient needed to discharge (S) and complete the proof of $(\star)$.
\end{enumerate}

\section{Summary}

We have proved, in full and rigorously: the closed form for $\Phi_n$
(Lemma~\ref{lem:phiform}); covariance and the identity element
(Lemma~\ref{lem:cov}); the Gaussian convolution semigroup and the fact that
finite free convolution with $G_t$ is the heat flow
$\partial_t p=-\tfrac12 p''$ (Lemma~\ref{lem:gauss},
Proposition~\ref{prop:heatconv}); the identification of root velocity with the
score (Lemma~\ref{lem:rootvel}); the exact de Bruijn identity with nonnegative
dissipation, hence monotonicity of $\Phi_n$ along the heat flow
(Theorem~\ref{thm:debruijn}); the full result with equality for $n=2$
(Proposition~\ref{prop:n2}); and the reduction of the general inequality
$(\star)$ to a single infinitesimal Stam inequality (S) along a jointly
coupled heat flow (Theorem~\ref{thm:reduction}), together with the precise
identification of the remaining ingredient (a finite free score-projection
lemma) and numerical evidence for it.

\end{document}
